\begin{figure}[htbp]
\centering
\begin{tikzpicture}
\begin{axis}[
  xbar, bar width=10pt, width=0.8\textwidth, height=6.0cm,
  xmin=0, xmax=24, xlabel={IPC gain of the V-Cache CCD over the plain CCD (\%)},
  symbolic y coords={libx264,GHPC,7-Zip,Cyberpunk 2077,Call of Duty},
  ytick=data, y dir=reverse, yticklabel style={font=\small},
  nodes near coords, nodes near coords style={font=\small,anchor=west},
  point meta=explicit symbolic,
  xmajorgrids, grid style={draftgray!25}, axis y line*=left, axis x line*=bottom,
  enlarge y limits=0.15]
\addplot[fill=draftlight,draw=draftblue] coordinates {
  (4.9,libx264)[4.9]
  (9.67,GHPC)[9.7]
  (9.75,7-Zip)[9.8]
  (13.5,Cyberpunk 2077)[13.5]
  (19.7,Call of Duty)[19.7]};
\end{axis}
\end{tikzpicture}
\caption[IPC gain from 3D V-Cache with clock, core design and memory held equal]{IPC gain from 3D V-Cache with clock, core design and memory held equal: the two CCDs of one Ryzen~9 7950X3D, both fixed at 4.2\,GHz, measured with performance counters by Lam \cite{lamAMDs7950X3DZen2023}. Cyberpunk~2077 is computed from the reported IPC of 1.26 and 1.43; the Call of Duty value is computed from the IPC of 0.71 and 0.85 read from Lam's chart (the text says ``over 19\%''). A sixth workload, a flight simulator with very few L3 misses, lost a small, unquantified amount of IPC and is not plotted.}
\label{fig:lam-ipc}
\end{figure}
